%%%%%%%%%%%%%%%%%%%%%%%%%%%%%%%%%%%%%%%%%%%%%%%%%%%%%%%%%%%%%%%%%%%%%%
% Lớp 10 | Bất phương trình bậc 2 một ẩn | Giải bất phương trình bậc 2 một ẩn
% Dạng 0D7?2-5: Bất phương trình chứa căn, | · |
% Số câu: 2  (sinh tự động bởi tach_cau.py; từ nay thầy sửa trực tiếp file này)
%%%%%%%%%%%%%%%%%%%%%%%%%%%%%%%%%%%%%%%%%%%%%%%%%%%%%%%%%%%%%%%%%%%%%%

%%%%%%%%%%%%%%%% TRẢ LỜI NGẮN %%%%%%%%%%%%%%%%

% nguồn: MaHoaNH_TraLoiNgan_Lop10
\begin{bt}%[0D7H2-5]%[ma:118b838e]
	Gọi $S$ là tập nghiệm của phương  trình $\sqrt{2x^2-x+3}=x+1$. Tổng lập phương tất cả các phần tử của $S$ bằng bao nhiêu?
	 \shortans{9}
	\loigiai{ 
		Ta có $\sqrt{2x^2-x+3}=x+1\Rightarrow  2x^2-x+3=(x+1)^2\Rightarrow x^2-3x+2=0\Rightarrow \hoac{&x=1\\&x=2.}$\\
		 Thử lại ta thấy hai nghiệm này đều thỏa mãn phương trình đã cho.\\
		Vậy nghiệm của phương trình đã  cho là $x=1$, $x=2$ và $1^3+2^3=9$.
	} 
\end{bt}

% nguồn: MaHoaNH_TraLoiNgan_Lop10
\begin{bt}%[0D7H2-5]%[ma:8426019f]
	Số nghiệm  của phương trình sau $\sqrt{2x^2+5x+3}=-3-x$ là
	\shortans{0} 
	\loigiai{ 
		Bình  phương hai vế phương trình đã cho, ta được
		\begin{eqnarray*}
			 &&2x^2+5x+3=(-3-x)^2\\
			&\Rightarrow&  2x^2+5x+3=9+6x+x^2\\
			&\Rightarrow& x^2-x-6=0\\ 
			 &\Rightarrow& \hoac{&x=3\\&x=-2.}
		\end{eqnarray*}
		Thay lần lượt các giá trị trên vào phương trình đã cho ta thấy không có giá trị nào  thỏa mãn.\\
		Vậy phương trình  đã cho vô nghiệm.
	 }
\end{bt}
